\item[\textbf{Exercise 2.1-3}]
Rewrite the \textsc{Insertion-Sort} procedure to sort into monotonically
decreasing instead of monotonically increasing order.

\Solution

The original \textsc{Insertion-Sort} procedure from the textbook (Chapter 2,
Section 1, page 19) is as follows:

\begin{codebox}
\Procname{$\proc{Insertion-Sort}(A,n)$}
\li \For $i \gets 2$ \To $n$ \Do
\li     $\id{key} \gets A[i]$
\li     \Comment{Insert $A[i]$ into the sorted subarray $A[1 \subarr i-1]$.}
\li     $j \gets i - 1$
\li     \While $j > 0$ and $A[j] > \id{key}$ \Do
\li         $A[j + 1] \gets A[j]$
\li         $j \gets j - 1$
        \End
\li     $A[j + 1] \gets \id{key}$
    \End
\end{codebox}

The comparison logic governing the sort order occurs on line 5. The condition
$A[j] > \id{key}$ results in an ascending (increasing) order because it shifts
elements to the right only if they are greater than the key. To sort in
monotonically decreasing order, we reverse the comparison operator to
$A[j] < \id{key}$.

The modified code with the change highlighted in red is:

\begin{codebox}
\Procname{$\proc{Insertion-Sort}(A,n)$}
\li \For $i \gets 2$ \To $n$ \Do
\li     $\id{key} \gets A[i]$
\li     \Comment{Insert $A[i]$ into the sorted subarray $A[1 \subarr i-1]$.}
\li     $j \gets i - 1$
\li     \While $j > 0$ and $\textcolor{red}{A[j] < \id{key}}$ \Do
\li         $A[j + 1] \gets A[j]$
\li         $j \gets j - 1$
        \End
\li     $A[j + 1] \gets \id{key}$
    \End
\end{codebox}
